\section{A pressure law with a quadratic budget}
\label{pressure:section}

The local metric deformation will use a pressure depending both on a height
variable $x$ and on the drift parameter $t$.  Its derivative in $x$ must be
large, while a large positive derivative in $x$ must be accompanied by a
favorable derivative in $t$.  The following construction supplies these
properties with constants chosen before any local geometric data.

\begin{theorem}[Pressure law]\label{pressure:law}
For every $A_0,P_0,\eps,b_0>0$ there exist $0<S_0\le1$, $T\ge1$, and a
smooth function
\[
 F:(-S_0,S_0)\times[T,\infty)\longrightarrow\R
\]
which is odd in $x$, has $F_x>0$, and satisfies
\begin{equation}\label{pressure:estimates}
\begin{gathered}
 F_x\ge A_0t,\qquad |F_t|\le\eps|F|,\qquad FF_t\le\eps F_x,\\
 F_x\le b_0t^2\quad\text{or}\quad -FF_t\ge P_0F_x,\\
 F_x\le b_0t^2+\eps F^2.
\end{gathered}
\end{equation}
For each fixed $t\ge T$,
\[
 \lim_{x\to-S_0}F(x,t)=-\infty,
 \qquad \lim_{x\to S_0}F(x,t)=+\infty.
\]
These limits are uniform when $t$ ranges over a compact subinterval of
$[T,\infty)$.  The choices depend only on $A_0,P_0,\eps,b_0$ and fixed
numerical constants.  In particular, they are independent of any manifold,
metric, local cutoff, or subsequent deformation parameter.  The lower bound
$T$ may be increased after the construction.
\end{theorem}

\begin{proof}
We construct $F$ by prescribing its inverse in the first variable. For
the moment, let $k>0$ and let $g$ be a positive smooth function on
$[0,\infty)$, both to be chosen. Write $f>0$ for an independent pressure
value and, for $t>0$, set
\begin{equation}\label{pressure:inverse}
 X(f,t)=\frac1k\arctan\frac{f^{1/3}}{g(y)},\qquad
 q=f^{2/3},\qquad y=\frac tq.
\end{equation}
Put $A(y)=g(y)+2yg'(y)$. Differentiation at fixed $f$ or at fixed $t$,
respectively, gives
\begin{equation}\label{pressure:inverse-derivatives}
 fX_t=-\frac{g'(y)}{k(1+g(y)^2/q)},\qquad
 X_f=\frac{A(y)}{3kq^2(1+g(y)^2/q)}.
\end{equation}
Thus $A>0$ will make $X$ invertible in $f$. For the inverse,
$F_x=1/X_f$ and $-FF_t/F_x=fX_t$. We will arrange the favorable
pressure branch on a decreasing part of $g$, and obtain the quadratic
branch from a lower bound for $t^2X_f$. We choose $g$ to decrease for
small $y$ and to equal $y$ for large $y$, while keeping $A$ positive
through the transition.

We construct $g$ as a smooth rounded maximum with a controlled negative
derivative.  Fix a nonnegative smooth probability density $\rho$ supported
in $[1,2]$, positive in $(1,2)$, and vanishing to infinite order at its
endpoints.  For example, normalize
\[
 \rho(r)=
 \begin{cases}
 c_\rho\exp\bigl(-1/((r-1)(2-r))\bigr),&1<r<2,\\
 0,&r\notin(1,2),
 \end{cases}
 \qquad \int_1^2\rho(r)\,\dd r=1.
\]
Write $\mu=\int_1^2r\rho(r)\,\dd r$, so $1<\mu<2$.  For $D_g>0$ set
\begin{equation}\label{pressure:rounded-max}
 h(y)=D_g(1+y)^{-1/4},\qquad
 g(y)=\int_1^2\max\{y,rh(y)\}\rho(r)\,\dd r,
 \qquad y\ge0.
\end{equation}
Here and below smoothness at $y=0$ means smooth extendibility to a
neighborhood of $[0,\infty)$.  To verify smoothness, put
\[
 M(z)=\int_1^2\max\{z,r\}\rho(r)\,\dd r.
\]
Then $M''=\rho$, with $\rho$ extended by zero, and
$g(y)=h(y)M(y/h(y))$.  Thus $g$ is smooth; moreover $M(z)=\mu$ for
$z\le1$ and $M(z)=z$ for $z\ge2$.  Consequently $g>0$, $g\ge y$, and
$g(y)=y$ for all sufficiently large $y$.

For a direct derivative estimate, let
\[
 L(z)=\int_{r\le z}\rho(r)\,\dd r,
 \qquad R(z)=\int_{r\ge z}r\rho(r)\,\dd r.
\]
At $z=y/h(y)$ the identities
\[
 g=yL(z)+hR(z),\qquad
 g'=L(z)-\frac{hR(z)}{4(1+y)}
\]
give
\begin{equation}\label{pressure:g-bounds}
 g\ge y,\qquad
 -\frac{g}{4(1+y)}\le g'\le1,\qquad
 A(y)=g(y)+2yg'(y)\ge\frac{y+2}{2(y+1)}g(y)\ge\frac12g(y).
\end{equation}
In particular the lower bound for $A/g$ is numerical and does not
deteriorate as $D_g$ increases.

We now specify the order of all choices.  Choose
\begin{equation}\label{pressure:k-Y}
 k\ge\max\{A_0,\eps^{-1},\pi/2\},\qquad
 Y\ge1\quad\text{with}\quad Y^{-3}\le\frac{b_0}{12k}.
\end{equation}
Choose $D_g$ only after $k,Y$, large enough that
\begin{equation}\label{pressure:D-choice}
 D_g\ge\max\left\{
 Y(1+Y)^{1/4},\,
 \frac{8kP_0(1+Y)^{5/4}}{\mu}
 \right\}.
\end{equation}
Since $y(1+y)^{1/4}$ is increasing, $y/h(y)\le1$ on $[0,Y]$.
Hence throughout this interval
\begin{equation}\label{pressure:decreasing-range}
 g(y)=\mu D_g(1+y)^{-1/4},\qquad
 g'(y)=-\frac{\mu D_g}{4(1+y)^{5/4}}\le-2kP_0.
\end{equation}
The equality cases at the endpoints are covered by the smooth matching
already proved.

With this $g$ fixed, define the finite constants
\[
 g_0=\min_{0\le y\le Y}g(y)>0,\qquad
 G_0=\max_{0\le y\le Y}g(y),\qquad
 B=\sup_{y>Y}\frac{g(y)}y<\infty.
\]
For example,
\[
 1\le B\le\max\left\{1,\frac{2D_g}{Y(1+Y)^{1/4}}\right\},
\]
because $g(y)\le\max\{y,2D_g(1+y)^{-1/4}\}$.
Finally choose
\begin{equation}\label{pressure:T-choice}
 T\ge\max\left\{
 1,\frac6\eps,\frac{12kB^2}{b_0},\,
 YG_0^2,\frac{12kY}{\eps g_0}
 \right\}.
\end{equation}
Thus the order is $k$, then $Y$, then $D_g$ and $g$, and then $T$;
none of the preceding choices depends on a later one.

We now invert the map $X$ in \cref{pressure:inverse}.
By \cref{pressure:g-bounds}, $X_f>0$.  As $f\downarrow0$, the argument
$y$ tends to infinity, where $g(y)=y$, and $X(f,t)\to0$.  As
$f\to\infty$, $y\to0$, $g(y)\to g(0)>0$, and
$X(f,t)\to S_0:=\pi/(2k)\le1$.  Thus for every $t>0$, $X(\cdot,t)$
maps $(0,\infty)$ bijectively onto $(0,S_0)$ with positive derivative.
Its inverse $f=F(x,t)$ is smooth by the implicit function theorem.

On a neighborhood of $f=0$ the exact identity $g(y)=y$ gives
\[
 X(f,t)=\frac1k\arctan(f/t),\qquad F(x,t)=t\tan(kx).
\]
Define $F(0,t)=0$ and extend $F$ oddly to $-S_0<x<0$.  The displayed
formula proves joint smoothness across $x=0$, including at the endpoint
$t=T$ by smooth extension to $t>0$.  The inverse construction also
proves the asserted pole limits.  For uniformity, if $I\subset[T,\infty)$
is compact and $H>0$, then $\max_{t\in I}X(H,t)<S_0$.  Whenever $x$
exceeds this maximum, monotonicity gives $F(x,t)>H$ for every $t\in I$.
The negative pole follows by oddness.

It remains to prove \cref{pressure:estimates}.  For $x>0$ set
$f=F(x,t)$ and use the variables in \cref{pressure:inverse}.  Implicit
differentiation of $X(F(x,t),t)=x$ yields
\begin{equation}\label{pressure:signed-identities}
 F_x=\frac1{X_f},\qquad F_t=-\frac{X_t}{X_f},\qquad
 \frac{FF_t}{F_x}=\frac{g'}{k(1+g^2/q)},\qquad
 \frac{F_t}{F}=\frac{3yg'}{t(g+2yg')}.
\end{equation}
In particular the sign in the last identity is the sign of $g'$.
Since $A\le g+2y\le3g$ and $y\le g$, we have
\begin{equation}\label{pressure:lower-slope}
 tX_f=\frac{yA}{3k(q+g^2)}\le\frac1k.
\end{equation}
This proves $F_x\ge kt\ge A_0t$ without a constant depending on $g$.
Also $|yg'|\le g$: for $g'\ge0$ this follows from $g'\le1$ and
$y\le g$, whereas for $g'<0$ the lower derivative bound in
\cref{pressure:g-bounds} even gives $|yg'|\le g/4$.  Therefore
\[
 \left|\frac{F_t}{F}\right|\le\frac6t\le\eps,
 \qquad
 \frac{FF_t}{F_x}\le\frac1k\le\eps.
\]
For the second inequality a negative numerator is already favorable;
for a nonnegative numerator use $g'\le1$.

Consider first $y>Y$.  Using $A\ge g/2\ge y/2$ and $g/y\le B$ gives
\begin{align}
 t^2X_f
 &=\frac{y^2A}{3k(1+yg^2/t)}\notag\\
 &\ge\frac{y^3}{6k(1+B^2y^3/T)}
  =\frac1{6k(y^{-3}+B^2/T)}
  \ge\frac1{6k(Y^{-3}+B^2/T)}
  \ge\frac1{b_0}.
 \label{pressure:large-y}
\end{align}
The last inequality uses \cref{pressure:k-Y,pressure:T-choice}.
This bound remains uniform as $y\to\infty$, and it proves
$F_x\le b_0t^2$ throughout this region.

If $0<y\le Y$, then $q=t/y\ge T/Y$ and
\[
 \frac{g^2}{q}\le\frac{YG_0^2}{T}\le1.
\]
It follows from \cref{pressure:decreasing-range,pressure:inverse-derivatives}
that
\begin{equation}\label{pressure:favorable-branch}
 -\frac{FF_t}{F_x}=fX_t
   =\frac{-g'}{k(1+g^2/q)}\ge P_0.
\end{equation}
The two regions therefore establish the alternative in
\cref{pressure:estimates}.  In this second region we also have
\begin{equation}\label{pressure:quadratic-upper}
 F^2X_f=\frac{qA}{3k(1+g^2/q)}
   \ge\frac{qg_0}{12k}
   \ge\frac{g_0T}{12kY}\ge\frac1\eps,
\end{equation}
so $F_x\le\eps F^2$.  Together with the small-$F_x$ bound for $y>Y$
this proves the final inequality of \cref{pressure:estimates}.

Oddness preserves all five inequalities on $x<0$, since $F_x$, $F^2$,
and $FF_t$ are even and $|F_t|,|F|$ are even.  At $x=0$ they follow
by continuity; more explicitly, $F=F_t=0$ and $F_x=kt$, while
$T\ge12kB^2/b_0\ge k/b_0$ gives the small-$F_x$ branch there.
All claims follow.
\end{proof}

\begin{remark}\label{pressure:parameter-independence}
When $A_0,P_0,\eps,b_0$ are prescribed using only the dimension and fixed
numerical constants, the same is true of $S_0,T$ and the entire pressure
law.  In particular, the small parameters denoted by $a$ and $d$ in the
local deformation may be chosen after this theorem has been applied,
without modifying the pressure or increasing $T$.
\end{remark}
